% from ega0-1-fr.tex
\oldpage[0\textsubscript{I}]{46}

De même, si $u:\mathcal{F}\to\mathcal{G}$ est un homomorphisme tel que
$u_x=0$, alors il existe un voisinage $U$ de $x$ tel que $u_y=0$ pour tout
$y\in U$.
\end{env}

\begin{env}[5.2.3]
\label{0.5.2.3-fr}
Supposons $X$ \emph{quasi-compact}, et soient $\mathcal{F}$, $\mathcal{G}$ deux
$\mathcal{O}_X$-Modules tels que $\mathcal{G}$ soit \emph{de type fini},
$u:\mathcal{F}\to\mathcal{G}$ un homomorphisme \emph{surjectif}. Supposons de
plus que $\mathcal{F}$ soit limite inductive d'un système inductif
$(\mathcal{F}_\lambda)$ de $\mathcal{O}_X$-Modules. Alors il existe un indice
$\mu$ tel que l'homomorphisme $\mathcal{F}_\mu\to\mathcal{G}$ soit
\emph{surjectif}. En effet, pour tout $x\in X$, il existe un système fini de
sections $s_i$ de $\mathcal{G}$ au-dessus d'un voisinage ouvert $U(x)$ de $x$
tel que les $(s_i)_y$ engendrent $\mathcal{G}_y$ pour tout $y\in U(x)$~; il y
a donc un voisinage ouvert $V(x)\subset U(x)$ de $x$ et $n$ sections $t_i$ de
$\mathcal{F}$ au-dessus de $V(x)$ telles que
$s_i|V(x)=u(t_i)$ pour tout $i$~; on peut en outre supposer que les $t_i$ sont
les images canoniques de sections d'un même faisceau
$\mathcal{F}_{\lambda(x)}$ au-dessus de $V(x)$. Couvrons alors $X$ avec un
nombre fini de voisinages $V(x_k)$, et soit $\mu$ un indice supérieur à tous
les $\lambda(x_k)$~; il est clair que cet indice répond à la question.

Supposant toujours $X$ quasi-compact, soit $\mathcal{F}$ un
$\mathcal{O}_X$-Module de type fini engendré par ses sections au-dessus de $X$
(5.1.1)~; alors $\mathcal{F}$ est engendré par une sous-famille \emph{finie}
de ces sections~: il suffit en effet de couvrir $X$ par un nombre fini de
voisinages ouverts $U_k$ tels que, pour chaque $k$, il y ait un nombre fini de
sections $s_{ik}$ de $\mathcal{F}$ au-dessus de $X$ dont les restrictions à
$U_k$ engendrent $\mathcal{F}|U_k$~; il est clair que les $s_{ik}$ engendrent
alors $\mathcal{F}$.
\end{env}

\begin{env}[5.2.4]
\label{0.5.2.4-fr}
Soit $f:X\to Y$ un morphisme d'espaces annelés. Si $\mathcal{G}$ est un
$\mathcal{O}_Y$-Module de type fini, alors $f^*(\mathcal{G})$ est un
$\mathcal{O}_X$-module de type fini. En effet, pour tout $x\in X$, il y a un
voisinage ouvert $V$ de $f(x)$ dans $Y$ et un homomorphisme surjectif
$v:\mathcal{O}_Y^p|V\to\mathcal{G}|V$. Si $U=f^{-1}(V)$ et si $f_U$ est la
restriction de $f$ à $U$, on a
$f^*(\mathcal{G})|U=f_U^*(\mathcal{G}|V)$~; comme $f_U$ est exact à droite
(4.3.1) et commute aux sommes directes (4.3.2), $f_U^*(v)$ est un
homomorphisme surjectif
$\mathcal{O}_X^p|U\to f^*(\mathcal{G})|U$.
\end{env}

\begin{env}[5.2.5]
\label{0.5.2.5-fr}
On dit qu'un $\mathcal{O}_X$-Module $\mathcal{F}$ \emph{admet une présentation
finie} si, pour tout $x\in X$, il existe un voisinage ouvert $U$ de $x$ tel
que $\mathcal{F}|U$ soit isomorphe à un \emph{conoyau d'un
$(\mathcal{O}_X|U)$-homomorphisme}
$\mathcal{O}_X^p|U\to\mathcal{O}_X^q|U$, $p$ et $q$ étant deux entiers $>0$.
Un tel $\mathcal{O}_X$-Module est donc de type fini et quasi-cohérent. Si
$f:X\to Y$ est un morphisme d'espaces annelés, et si $\mathcal{G}$ est un
$\mathcal{O}_Y$-Module admettant une présentation finie,
$f^*(\mathcal{G})$ admet une présentation finie, comme le montre le
raisonnement de (5.1.4).
\end{env}

\begin{env}[5.2.6]
\label{0.5.2.6-fr}
Soit $\mathcal{F}$ un $\mathcal{O}_X$-Module admettant une présentation finie
(5.2.5)~; alors, pour tout $\mathcal{O}_X$-Module $\mathcal{H}$,
l'homomorphisme canonique fonctoriel
\[
  \bigl(\mathcal{H}\!om_{\mathcal{O}_X}(\mathcal{F},\mathcal{H})\bigr)_x
  \longrightarrow
  \operatorname{Hom}_{\mathcal{O}_x}(\mathcal{F}_x,\mathcal{H}_x)
\]
est \emph{bijectif} (T, 4.1.1).
\end{env}

\begin{env}[5.2.7]
\label{0.5.2.7-fr}
Soient $\mathcal{F}$, $\mathcal{G}$ deux $\mathcal{O}_X$-Modules ayant une
présentation finie. Si, pour un $x\in X$, $\mathcal{F}_x$ et $\mathcal{G}_x$
sont deux $\mathcal{O}_x$-modules \emph{isomorphes}, alors il existe un
voisinage ouvert $U$ de $x$ tel que $\mathcal{F}|U$ et $\mathcal{G}|U$ soient
\emph{isomorphes}. En effet, si
$\varphi:\mathcal{F}_x\to\mathcal{G}_x$ et
$\psi:\mathcal{G}_x\to\mathcal{F}_x$ sont deux isomorphismes réciproques, il
existe, d'après (5.2.6), un voisinage ouvert $V$ de $x$ et une section $u$
(resp. $v$) de
$\mathcal{H}\!om_{\mathcal{O}_X}(\mathcal{F},\mathcal{G})$ (resp.
$\mathcal{H}\!om_{\mathcal{O}_X}(\mathcal{G},\mathcal{F})$) au-dessus de $V$
telle

